\documentclass[11pt,a4paper]{article}
\usepackage[T1]{fontenc}
\usepackage{lmodern,amsmath,amssymb,booktabs,microtype}
\usepackage[margin=24mm]{geometry}
\usepackage[colorlinks=true,urlcolor=blue,linkcolor=blue]{hyperref}
\newcommand{\EE}{\mathbb E}
\newcommand{\tr}{\operatorname{tr}}
\newcommand{\W}{\mathsf W}
\title{Uniform rank bounds for finite-time propagation\\From independent products to all coupling chaoses}
\author{Alan Day \and Chengheng\thanks{AI research assistant (OpenAI).}}
\date{Local research note --- 7 November 2026}
\begin{document}\maketitle
\begin{abstract}
We turn the established rank-product covariance formula into a
positive Markov chain. Its first two conditional Casimir moments
give a rank-deficit bound uniform in the number of factors, not
merely at fixed polynomial degree. Gaussian tensor symmetrization
transfers the bound to the highest coupling chaos. A one-step
version and Gaussian creation/annihilation operators then control
the rank--chaos deficit of the full ordinary propagator, including
all lower contractions, uniformly on compact real-time intervals.
A finite-parameter covariance-filter estimate follows. We also identify an exact
charged-probe four-point interface and a positive regression
remainder sufficient for the actual memory-field bound. Connecting
the rank--chaos estimate to the charged force remains the next step.
\end{abstract}
\section{State, model and the all-degree statement}
Use standard CAR, $\hbar=1$, string time, and the overlapping
complex Gaussian model $H=\sum_{I,J}J_{IJ}c_I^\dagger c_J$, with
\begin{equation}
 \EE J_{IJ}J_{KL}=w_N\delta_{IL}\delta_{JK},\quad
 w_N=4^p\frac{J_s^2}{\lambda_N}\left(\frac{p!}{N^p}\right)^2,
 \qquad \lambda_N=p^2/N.
\end{equation}
$J=J^\dagger$; all charge sectors use the same coupling.
At charge $m$, $\tr_m=\operatorname{Tr}/\binom Nm$ and
$\|X\|_{2,m}^2=\EE\tr_mX^\dagger X$. Put
\begin{equation}
 C=C_m=w_N\binom mp\binom{N-m+p}p,\qquad
 \Phi(X)=\EE_{\widetilde H}\widetilde H X\widetilde H.
\end{equation}
Here $\widetilde H$ is an independent copy; $\Phi$ is a deterministic linear map.
The neutral operator algebra has multiplicity-free $U(N)$ ranks
$0\le r\le R_m=\min(m,N-m)$, with orthogonal projectors $P_r$.
These are representation ranks, not matrix ranks.

For independent copies $H_1,\ldots,H_k$ in this sector define
\begin{equation}
 \pi_k(r)=C^{-k}\EE\tr_m|P_r(H_1\cdots H_k)|^2,\qquad D_k(r)=kp-r.
\end{equation}
For $N\ge8$, $1\le p\le N/8$, $N/4\le m\le3N/4$, and every
integer $k\ge1$, we prove
\begin{equation}
 \boxed{\ \EE_{\pi_k}D_k^2
 \le A\big(\lambda_N k^2+\lambda_N^2k^4\big),\qquad A=2\cdot10^6.\ }
 \label{eq:main}
\end{equation}
The constant is deliberately unoptimized. The result is a finite
parameter bound. For double scaling take even $N\to\infty$,
$\lambda_N\to\lambda>0$ in a positive compact interval. The full
$\beta=\mu=0$ state uses $b_m=2^{-N}\binom Nm$.

\newpage
\section{An exact rank-convolution Markov chain}
The existing covariance and Racah interfaces \cite{Kota,DLMF} give
\begin{align}
 \Phi(P_sX)&=\kappa_sP_sX,&
 \kappa_s&=w_N\binom{m-s}p\binom{N-m+p-s}p,\\
 \chi_r(s)&={}_4F_3\!\left(\begin{matrix}-r,r-N-1,-s,s-N-1\\
                  -m,m-N,-N-1\end{matrix};1\right),&
 \mu_s&=\kappa_s/C.
\end{align}
Invalid binomials are zero. Let $d=\binom Nm$ and
$d_r=\binom Nr^2-\binom N{r-1}^2$. Covariance composition gives
\begin{equation}
 \sum_r\pi_k(r)\chi_r(s)=\mu_s^k,\qquad
 \pi_k(r)=\frac{d_r}{d^2}\sum_s d_s\chi_r(s)\mu_s^k.
 \label{eq:transform}
\end{equation}
Positivity is not inferred from the signs of this sum. Decompose
$H=\sum_a H_a$ into independent isotropic rank components and let
$\omega_a=C^{-1}\EE\tr_mH_a^2$. If $\pi_{ra}(t)$ is the known
positive output law of two independent normalized pure-rank
inputs, the transition kernel is
\begin{equation}
 P(r,t)=\sum_a\omega_a\pi_{ra}(t),\qquad
 \pi_{k+1}(t)=\sum_r\pi_k(r)P(r,t),\qquad \pi_0=\delta_0.
\end{equation}
Indeed the sesquilinear coefficient covariance of $H_1\cdots H_k$
is $U(N)$ invariant and therefore scalar on each neutral rank.
No Hermiticity of that product is required. Its covariance map is
$\Phi^k$, and the invariant Choi decomposition gives
\eqref{eq:transform}. Multiplication implies $r\le kp$ on its
support, hence $D_k\ge0$ even when $kp>R_m$.

Write $x_r=r(N+1-r)$, $X_j=x_{r_j}$ and $n=N-m$. Define
\begin{align}
 e&=mn(N+1)/N,&u&=\mu_1=\frac{(m-p)n}{m(n+p)},\\
 v&=\mu_2=u\frac{(m-p-1)(n-1)}{(m-1)(n+p-1)},&
 c&=e(1-u)=\frac{pn(N+1)}{n+p},\\
 b&=\frac{(N-1)(N-2)}{m(m-1)n(n-1)N(N+1)},&
 d_2&=\frac{2(N-1)}{Ne}+bN.
\end{align}
The first two characters are $\chi_r(1)=1-x_r/e$ and
$\chi_r(2)=1-d_2x_r+bx_r^2$. Applying the transition to them gives
\begin{align}
 \EE[X_{j+1}\mid X_j=x]&=ux+c,\label{eq:meanstep}\\
 \operatorname{Var}(X_{j+1}\mid X_j=x)
   &=(v-u^2)x^2+Bx+V_1,\label{eq:varstep}\\
 B&=\frac{d_2}{b}(u-v)-2uc,&
 V_1&=\frac{p^2n^2(N+1)(m-p+1)}{(n+p)^2(n+p-1)}.
\end{align}
In particular $\EE X_k=e(1-u^k)$ and
$\EE X_k^2=(v^k-1+d_2\EE X_k)/b$, exactly at finite $N$.

\newpage
\section{Uniform control without large-order cancellations}
The sign $0\le v\le u^2\le1$ follows by comparing the two factors
in $v/u$ to those in $u$. Direct simplification also gives
\begin{equation}
 B=\frac{up\{n(m^2+n^2-2)-p(m^2+2mn-n^2-2n)\}}
            {(m-1)(n+p)(n+p-1)}.
\end{equation}
Under the central-sector assumptions of \eqref{eq:main}, bound
the numerator termwise and use $m-1,n+p-1\ge N/8$. For example,
the following conservative bounds all hold:
\begin{equation}
 |B|\le102400p,\quad V_1\le512p^2N,\quad
 c\le8pN,\quad1-u\le16p/N.\label{eq:coarse}
\end{equation}
The total-variance recurrence, with $M_j=\EE X_j$, is
\begin{equation}
 \operatorname{Var}X_{j+1}
 =v\operatorname{Var}X_j+(v-u^2)M_j^2+BM_j+V_1.
\end{equation}
Since $M_j\le jc$, discarding the nonpositive quadratic term and
summing \eqref{eq:coarse} yields, for every $k$,
\begin{equation}
 \operatorname{Var}X_k\le410112\,k^2p^2N.\label{eq:variance}
\end{equation}
If $t=kp\le N/8$, use
\begin{align}
 x_t&=kx_p-k(k-1)p^2,&
 x_p-c&=\frac{p^2(m-p+1)}{n+p},\\
 kc-M_k&=c\sum_{j<k}(1-u^j)
 & &\le\tfrac12k(k-1)c(1-u).
\end{align}
Together with support $r\le t$, these imply
$0\le x_t-M_k\le73k^2p^2$.
Furthermore $x_t-x_r=(t-r)(N+1-t-r)\ge(N/2)(t-r)$.
Thus
\begin{equation}
 \EE D_k^2\le\frac4{N^2}\big[\operatorname{Var}X_k+(x_t-M_k)^2\big]
 \le1640448\lambda_Nk^2+21316\lambda_N^2k^4.
\end{equation}
For $kp>N/8$, do not use monotonicity of $x_{kp}$:
simply $D_k^2\le k^2p^2\le64\lambda_N^2k^4$.
This proves \eqref{eq:main} uniformly over all degrees.

At half filling and fixed $k$, the exact formulas also give
\begin{equation}
 \EE D_k\to\lambda k^2,\qquad
 \operatorname{Var}D_k\to\lambda k^2,\qquad
 \EE D_k^2\to\lambda^2k^4+\lambda k^2.\label{eq:DS}
\end{equation}
For the mean, expand the exact $e(1-u^k)$ at fixed $k$.
For variance, $B/p\to2$, $V_1/N^2\to\lambda$,
$(v-u^2)p^2\to0$, and $M_j/(Np)\to j$ turn the recurrence into
increments $\lambda(2j+1)$. Finally
$(x_{kp}-X_k)/N=D_k[1+O(kp/N)]$ uniformly on its support.
These two moments match a Poisson law, but do not prove that law.

\newpage
\section{Highest chaos, Wick series and a covariance-filter bound}
Now return to the same coupling in every factor. Let
$Z_k=\W_k(H^k)$, where $\W_k$ is the orthogonal projection onto
degree-$k$ Gaussian coupling chaos. In real Gaussian coordinates,
let $T$ be the ordered $k$-linear coefficient tensor of the product.
For any deterministic rank weight $w$, the Gaussian isometry gives
\begin{equation}
 \|wZ_k\|_{2,m}^2=k!\|\operatorname{Sym}(wT)\|^2
 \le k!\|wT\|^2=k!C^k\EE_{\pi_k}|w(r)|^2.\label{eq:chaos}
\end{equation}
Symmetrization is an orthogonal projection on the input tensor;
it commutes with the output rank filter. There is only one $k!$.

Let $\mathcal D_k=\sum_r(kp-r)P_r$. Distinct coupling chaoses
are orthogonal. Therefore the series
\begin{equation}
 \mathcal D(s)=\sum_{k\ge0}\frac{(is)^k}{k!}\mathcal D_k Z_k
\end{equation}
is entire in annealed $L^2$. With $x=C|s|^2$, its squared norm obeys
\begin{equation}
 \|\mathcal D(s)\|_{2,m}^2\le A e^x
 \left[\lambda_N(x^2+x)
 +\lambda_N^2(x^4+6x^3+7x^2+x)\right].\label{eq:wick}
\end{equation}
This is an all-degree sum, using the exact exponential generating
functions for $k^2$ and $k^4$, not a time truncation.
For full charge summation use the previously established
$C_m\le M$ uniformly in $m,N$. Outside the central region,
$D_k\le kp$ and \eqref{eq:chaos} give the total tail bound
\begin{equation}
 2p^2e^{-N/8}e^x(x^2+x),\qquad x=M|s|^2.
\end{equation}
It vanishes on compact time sets. Inactive sectors are zero.

In the same central sectors there is a finite-parameter filter consequence. Extend the
exact covariance eigenvalue to $z\ge0$ by
\begin{equation}
 \mu(z)=\prod_{j=0}^{p-1}
 \left(1-\frac z{m-j}\right)_+
 \left(1-\frac z{N-m+p-j}\right)_+.
\end{equation}
This agrees with $\mu_r$ at every allowed integer rank. The product
is globally Lipschitz with constant at most
$p/(m-p+1)+p/(N-m+1)\le12p/N$. Consequently
\begin{align}
 |p(\mu_r-\mu(kp))|&\le12\lambda_ND_k(r),\\
 \|p(\Phi/C-\mu(kp))Z_k\|_{2,m}^2
 &\le144\lambda_N^2 k!C^k\EE_{\pi_k}D_k^2.\label{eq:filter}
\end{align}
Its corresponding Wick series is bounded by $144\lambda_N^2$
times \eqref{eq:wick}. Keeping $\mu(kp)$ exact avoids amplifying
the finite-parameter error of an early replacement by $q^k$.
Neither $Z_k$ nor this comparison term is asserted to commute with $H$.

\newpage
\section{A one-step bound controls creation and annihilation}
The all-degree estimate can be extended beyond the Wick series.
For one rank transition $r\to t$ with $r\le kp$, put
$\delta=p+r-t$. Multiplication by a rank-$a$ tensor has
$|t-r|\le a\le p$, hence $0\le\delta\le2p$.
For completeness, the lower triangle follows by cyclic trace:
if $r>t+a$, then $\langle C_t,A_rB_a\rangle
=\langle C_tB_a^\dagger,A_r\rangle=0$ since the former product
has rank at most $t+a$. No product multiplicity assumption is needed.

If $(k+1)p\le N/8$, \eqref{eq:meanstep}--\eqref{eq:coarse} give
\begin{align}
 0\le x_{r+p}-\EE_r x_t&\le32(k+1)p^2,\\
 \operatorname{Var}_r x_t&\le205312(k+1)p^2N.
\end{align}
For example the mean difference equals
$(1-u)x_r+[p(N+1)-c]-2rp-p^2$, and
$p(N+1)-c=p^2(N+1)/(n+p)$.
Using $\delta\le2(x_{r+p}-x_t)/N$ gives a bound
$821248\lambda_N(k+1)+4096\lambda_N^2(k+1)^2$.
If $(k+1)p>N/8$, use $\delta^2\le4p^2\le256\lambda_N^2(k+1)^2$.
Thus, with $A_1=10^6$, uniformly in all $k\ge0$,
\begin{equation}
 \sup_{r\le kp}\EE_r\delta^2
 \le A_1[\lambda_N(k+1)+\lambda_N^2(k+1)^2].\label{eq:step}
\end{equation}

Represent $H=\sum_\alpha g_\alpha B_\alpha$ with real standard
Gaussian coordinates, Hermitian $B_\alpha$, and
$\sum_\alpha B_\alpha^2=CI$. Work in Gaussian $L^2$ tensored with
the normalized matrix Hilbert space, restricted to $U(N)$-equivariant
functions. Let $\mathcal N$ count Gaussian chaos and
$\mathcal R=\sum_r rP_r$ act on the matrix output. The operator
\begin{equation}
 \mathfrak D=p\mathcal N-\mathcal R
\end{equation}
is nonnegative here: degree-$k$ equivariant coefficients have output
rank at most $kp$. Coupling tensors are built from ranks at most $p$.
Left multiplication by $H$ is
$\mathcal T=\mathcal C+\mathcal C^\dagger$, with
$\mathcal C=\sum_\alpha a_\alpha^\dagger L_{B_\alpha}$.

The deterministic operator
$\sum_\alpha L_{B_\alpha}P_tL_{B_\alpha}$ is $U(N)$ equivariant,
so its cross terms between different input ranks vanish exactly;
on rank $r$ its scalar is $CP(r,t)$. Apply this to the unsymmetrized
creation tensor, then use contraction of input symmetrization:
\begin{equation}
 \|[\mathfrak D,\mathcal C]\psi_k\|^2
 \le(k+1)C\sup_{r\le kp}\EE_r\delta^2\,\|\psi_k\|^2.
\end{equation}
The equivariant subspace is reduced by creation, annihilation,
number and rank. Thus the annihilation block is minus the adjoint
of this block. Adding the two adjacent blocks gives
\begin{equation}
 \|[\mathfrak D,\mathcal T]\psi\|^2
 \le4A_1C(\lambda_N+\lambda_N^2)
       \|(\mathcal N+1)^{3/2}\psi\|^2.\label{eq:commbound}
\end{equation}

\newpage
\section{The full ordinary propagator, with every chaos retained}
Let $\mathcal F(s)=e^{isH}$ and $s\in\mathbb R$. Its Gaussian
derivative tensor has a dimension-free bound. For any ordered
Duhamel integrand
\begin{equation}
 W_{\boldsymbol\alpha}=U_0B_{\alpha_1}U_1\cdots B_{\alpha_\ell}U_\ell,
 \qquad U_j\text{ unitary},
\end{equation}
the exact identity
$\sum_{\boldsymbol\alpha}\tr_m W_{\boldsymbol\alpha}^\dagger
W_{\boldsymbol\alpha}=C^\ell$ follows by successive summation
using $\sum B_\alpha^2=CI$ and cyclic trace. The $U_j=e^{it_jH}$
are allowed to depend on the sample. They do not depend on the
summed derivative indices.

Minkowski in the derivative-index $\ell^2$ and matrix Hilbert
space, followed by the $\ell!$ permutations and simplex integration,
therefore gives, already samplewise,
\begin{equation}
 \sum_{\boldsymbol\alpha}\tr_m
 |\partial_{\boldsymbol\alpha}\mathcal F(s)|^2
 \le C^\ell|s|^{2\ell}.\label{eq:derivatives}
\end{equation}
The factorial cancels the simplex volume. Gaussian chaos isometry
then implies, with $z=Cs^2$ and $\mathcal F_k=\W_k\mathcal F$,
\begin{align}
 \sum_k(k)_\ell\|\mathcal F_k(s)\|^2&\le z^\ell,\\
 \|(\mathcal N+1)^{3/2}\mathcal F(s)\|^2
 &\le1+7z+6z^2+z^3.\label{eq:number}
\end{align}
Here $(k)_\ell=k(k-1)\cdots(k-\ell+1)$. In fact summing the
nonnegative factorial moments yields
$\sum_k\zeta^k\|\mathcal F_k(s)\|^2\le e^{z(\zeta-1)}$ for
$\zeta\ge1$. This is a bound on the complete chaos distribution.

Since $\mathfrak D I=0$, the commutator Duhamel identity is
\begin{equation}
 \mathfrak D\mathcal F(s)
 =i\int_0^s e^{i(s-t)\mathcal T}
       [\mathfrak D,\mathcal T]\mathcal F(t)\,dt.
\end{equation}
At finite $N$, all Gaussian derivatives of $\mathcal F$ are bounded;
the number operator adds only polynomial growth in Gaussian coordinates.
The identity follows on the common Hermite core and its Gaussian
Sobolev closure. The left evolution is unitary for real time.
Equations \eqref{eq:commbound}--\eqref{eq:number} prove
\begin{equation}
 \boxed{\ \|\mathfrak D e^{isH}\|_{2,m}^2
 \le4A_1C(\lambda_N+\lambda_N^2)s^2(1+7z+6z^2+z^3).\ }
 \label{eq:ordinary}
\end{equation}
For full charge summation, $0\le\mathfrak D\le p\mathcal N$
on equivariant functions bounds the exterior-sector contribution by
$2p^2e^{-N/8}(z^2+z)$, now with $z=Ms^2$. Thus the actual ordinary
propagator has a uniformly bounded rank--chaos deficit on compact
real-time intervals, including all its lower contractions.

The Lipschitz argument also gives in central sectors
\begin{equation}
 \|p[\Phi/C-\mu(p\mathcal N)]e^{isH}\|_2
 \le12\lambda_N\|\mathfrak D e^{isH}\|_2.\label{eq:ordinaryfilter}
\end{equation}
The right comparison acts on Gaussian chaos; it is not asserted
to be a samplewise conserved function of $H$.

\newpage
\section{The actual finite-time target and a positive interface}
For the charge-removal dilation, let $h=H_m$, $K=1_N\otimes H_{m-1}$,
$V\psi=m^{-1/2}\sum_i|i\rangle c_i\psi$, $a=(1-p/m)h$, and
$R=KV-Va$. Write $F(s)=V^\dagger e^{iKs}V$ and
$\Lambda(s)=R^\dagger e^{iKs}R$.
The preceding stage showed that the actual pinched memory field
differs in $L^2$ by $o(1)$ from
\begin{equation}
 X_0(s)=(1-\mathsf P_h)m e^{-ihs/2}\Lambda(s)e^{-ihs/2},
 \qquad \mathsf P_hX=\sum_eP_eXP_e.
\end{equation}
Its uniform second moment is still the target.

Independent circular Gaussian vectors $\xi,\eta$, independent of
$J$, give linear charged probes $M_\xi=\sum_i\xi_i c_i$ with
$\EE_\xi\bar\xi_i\xi_j=\delta_{ij}$. Exactly at finite $N$,
\begin{equation}
 F(s)=\frac1m\EE_\xi M_\xi^\dagger e^{iH_{m-1}s}M_\xi.
\end{equation}
Two $F$ insertions are therefore a four-point function with two
independent probe colors. Define
\begin{equation}
 B_a(t_1,t_2)=e^{iat_1}F(t_2-t_1)e^{-iat_2}.
\end{equation}
Direct differentiation, using $\Lambda=-F''+i\{a,F'\}+aFa$, gives
\begin{equation}
 \partial_{t_1}\partial_{t_2}B_a(t_1,t_2)
 =e^{iat_1}\Lambda(t_2-t_1)e^{-iat_2}.\label{eq:probe}
\end{equation}
Thus four endpoint derivatives recover two ordinary force kernels
while retaining their common $J$. This does not model the forces
themselves as independent probes.

A positive subtraction avoids differentiating pinching. Define the
probability measure and its conditional function
\begin{align}
 \rho(dx,dy)&=\EE\tr_m[P_h(dx)V^\dagger P_K(dy)V],&
 \mu(dx)&=\rho(dx,\mathbb R),\\
 \ell_s(x)&=\EE_\rho[(y-(1-p/m)x)^2e^{isy}\mid x].
\end{align}
By the spectral expansion of $R$, $\ell_s(h)$ is the exact annealed
regression of $\Lambda(s)$ onto deterministic energy functions.
The Hilbert-space projection identity and pinching contraction give
\begin{align}
 \mathcal V_{N,m}(s)&=\EE\tr_m\Lambda^\dagger\Lambda
                  -\int|\ell_s|^2d\mu\ \ge0,\\
 \|X_0(s)\|_2^2&\le\sum_m b_m m^2\mathcal V_{N,m}(s).\label{eq:target}
\end{align}
The difference is residual sample-dependent conserved variance;
this bound is sufficient, not necessary. A uniform central-sector
estimate $\mathcal V_{N,m}=O(p^{-4})$, with the controlled tails,
would suffice. The established leading theorem yields only
$\sum_m b_m\mathcal V_{N,m}=o(p^{-2})$.
The rank--chaos estimate \eqref{eq:ordinary} does not by itself
bound this charged, energy-regressed quantity.

\newpage
\section{Literature boundary, verification and next calculation}
The finite-index trace rule in BNR \cite{BNR}, Eqs. (3.11)--(3.12),
retains all intersections. Its subsequent leading double-scale
simplification drops triple intersections. For a size-one charged
probe, one marked site lies in two independent Hamiltonian index
sets with probability $(p/N)^2=\lambda_N/N$; two probe indices
coincide with probability $1/N$. Both matter at connected order.
The published neutral-probe four-point formula cannot simply be
specialized to this charged, subleading problem. Endpoint
derivatives and regression subtraction must be controlled together.

The distinction that motivated Sections 5--6 is explicit:
\begin{equation}
 e^{isH}=\sum_{n\ge0}\frac{(is)^n}{n!}
          \sum_{k=0}^n\W_k(H^n),
 \qquad
 :e^{isH}:=\sum_{k\ge0}\frac{(is)^k}{k!}\W_k(H^k).
\end{equation}
The terms with $k<n$ contain lower-chaos contractions, already
included in \eqref{eq:ordinary}. For example,
$\W_1(H^3)=2CH+\Phi(H)$ need not be a scalar multiple of $H$;
a scalar Gaussian prefactor is not the argument used here.
The remaining task is the coherent transfer between adjacent
charge sectors in the actual force sandwich. It must connect
\eqref{eq:ordinaryfilter} to the charged four-point and the
energy regression \eqref{eq:target}; a Gaussian-chaos multiplier
cannot simply be declared to commute with $H$.
Every proposed limiting kernel must recover the earlier equal-time
dynamic variance
\begin{equation}
 W=4J_s^4\left[\frac1\lambda+\frac{4q}{1+q}\right],
 \qquad q=e^{-4\lambda}.
\end{equation}
An independent-site approximation that retains only the first term
misses the coherent common-coupling contribution.

\paragraph{Saved evidence.}
Exact rational Racah inversion and transition checks give 252
equalities over six finite systems and product lengths one through
eight. The rational bounds are checked in 1800 parameter instances;
four symbolic identities hold for general $N,p,m$. Six fixed-$k$
symbolic limits check \eqref{eq:DS}, whose general proof is the
recurrence above. Independently constructed CAR channels in three
systems reproduce the rank law through five factors. Nine complete
Gaussian coefficient tensors through degree three verify the
single-$k!$ bound. Eighteen complete Hermite coefficient cases
through $H^4$ check creation commutators including lower chaoses;
43 exact one-step laws check the conditional bound and rank triangle.
Three finite charged-probe systems verify the
four-point second-moment sum and endpoint derivative identity.
Five exact-formula evaluations at 80/120 digits agree within
$10^{-60}$; the large-$k$ case checks numerical evaluation, not a
fixed-$k$ asymptotic. These are reproducibility checks, not substitutes
for the uniform proof and not interval-certified rounding bounds.

\begin{thebibliography}{9}\small
\bibitem{Kota} V. K. B. Kota, \emph{Bivariate moments of the two-point
correlation function for embedded Gaussian unitary ensemble with
$k$-body interactions},
\href{https://arxiv.org/abs/2208.11312}{arXiv:2208.11312}, Sec. IV.
Used for covariance representation; spectral-density correlations
are distinct from the present operator covariance.
\bibitem{DLMF} NIST Digital Library of Mathematical Functions,
\href{https://dlmf.nist.gov/18.26.E3}{Eq. 18.26.3}, Racah polynomials.
The positive convolution and two-input interface were established
in the project's earlier rank-product note.
\bibitem{BNR} M. Berkooz, V. Narovlansky and H. Raj,
\emph{Complex Sachdev-Ye-Kitaev model in the double scaling limit},
\href{https://arxiv.org/abs/2006.13983}{arXiv:2006.13983},
Secs. 3.2--3.3, 7 and 8.1.
\end{thebibliography}
\end{document}
